% Physical unit 031; accepted Session 03 after fresh 400-dpi authority replay.
\GordanUnitStart{031}{printed page 26}
\GordanTopFolio{26}
\begingroup\small
Quam aequationem etsi longe elegantiorem formam ill. Jacobi ostendit induere posse, nos tamen ad rem nostram in illa stare possumus.

\sect{5}

Quod quia $\Theta$ functiones in series praecipiter convergentes explicantur, ideoque facillime computantur, jam in re erit, etiam functiones trigonometricas angulorum
\[
 B\qquad l\qquad T
\]
referre ad $\Theta$ functiones.

Formularum ea ratione evolutarum gravissimae hae sunt:
\[
\begin{gathered}
 \sin l=\sin l_0\sqrt{x'\over x}\,{H(u+K)\over\Theta(u)},\\
 \sqrt{x}\sin B=\sin B_0{H(u+K)\over\Theta(u+K)},\\
 {\cos B\over\cos l}=\sqrt{x'}{\Theta u\over\Theta(u+K)},\\
 \tg T=\cot l_0\sqrt{x'}\,{\Theta u\over H(u)}.
\end{gathered}
\]

\rule{0.35\textwidth}{0.4pt}

\caput{VII}
\sect{1}

Ut vero quantitates:
\[
 l\qquad s\qquad \psi
\]
ex aequationibus capitis superioris ad calculum vocentur, formulas statuere necesse est ad constantes:
\[
 l_0\qquad B_0\qquad \Sigma(l_0)\qquad \mathcal A(l_0)\qquad x\qquad x'\qquad K\qquad E\qquad q\qquad v
\]
omnes ex una computandas.

\sect{2}

Plerumque circulus, qui curva tangitur, definitus est; rarissimo ab hac ratione deceditur. Ubi altera e duabus periodis nota est; (e. g. cum quaeritur: \gq{Quemnam circulum tangit linea omnium brevissima duobus aequatoris punctis definita?}) tum ad computandum $l_0$ tantummodo seriebus Capit. III utemur.

\sect{3}

Perinde, $l_0$ nunc notum fingentes, determinare possumus quantitates:
\[
 B_0\qquad \Sigma(l_0)\qquad \mathcal A(l_0)\qquad x\qquad x'
\]
\endgroup
